%! Author = sriadityavedantam
%! Date = 3/29/26

Around this chapter is where most textbooks will start.
Some of the concepts that we introduced in the previous chapters are more seen as prior knowledge, or seen in an introduction chapter similar to what we have with Chapter~\ref{ch:chap1}.
So from here on out, consider everything you learned so far as a foundation for the rest of the content.

\section{Rings}

\begin{definition}[Ring]
    A ring $R$ is a set with two operations, $+$ and $\times$.
    Addition is commutative and associative.
    There exists an additive identity, usually denoted by $0_R$.
    There exists an additive inverse in $R$, say $b$, such that $a+b=0_R$, and this inverse is unique.
    Multiplication is associative.
    Together, the two operations satisfy the distributive laws.
    There may also be a multiplicative identity, denoted by $1_R$.
\end{definition}

\begin{definition}[Subrings]
    A subset $S$ is a subring of $R$ if for all $a,b\in S$, it has closure under addition and multiplication.
    It must also have the additive identity and additive inverses for each element.
\end{definition}

For example, in an introduction to proofs class we may have seen that
\[
    \mathbb{Z}\subseteq\mathbb{Q}\subseteq\mathbb{R}\subseteq\mathbb{C}.
\]
We learned them as sets, but looking at properties of rings and subrings, consider them all rings and subrings in that order.
However, if we wanted to look outside of these number systems, let us look at matrices:
\[
    \left\{
        \begin{bmatrix}
            a&b\\
            -b&a
        \end{bmatrix}
        :a,b\in\mathbb{Q}
    \right\}.
\]
Note that this is a subring of $\mat{2}(\mathbb{R})$.

\begin{definition}[Field]
    A field $\f$ is a commutative ring with $1\neq 0$ such that if $a\in\f$ and $a\neq 0$, then $a$ is a unit.
\end{definition}

\begin{definition}[Subfield]
    If $S$ is a subring of a field $\f$, and is also closed under multiplicative inverses of nonzero elements, then it is also a subfield.
\end{definition}

We have previously learned that
\[
    \left\{
        \begin{bmatrix}
            a&b\\
            -b&a
        \end{bmatrix}
        :a,b\in\mathbb{Q}
    \right\}
\]
is a subring of $\mat{2}(\mathbb{R})$.
But I also claim it is a field itself.

\begin{proposition}{
    The set
    \[
        \left\{
            \begin{bmatrix}
                a&b\\
                -b&a
            \end{bmatrix}
            :a,b\in\mathbb{Q}
        \right\}
    \]
    is a field.
}
    Suppose
    \[
        M=
        \begin{bmatrix}
            a&b\\
            -b&a
        \end{bmatrix}
        \neq
        \begin{bmatrix}
            0&0\\
            0&0
        \end{bmatrix},
    \]
    which means $a$ and $b$ are not both $0$.
    Then
    \[
        \det(M)=a^2+b^2.
    \]
    Since $a$ and $b$ are not both $0$, we have $a^2+b^2\neq 0$.
    Thus
    \[
        M^{-1}=\frac{1}{a^2+b^2}
        \begin{bmatrix}
            a&-b\\
            b&a
        \end{bmatrix}.
    \]
    Since $a,b\in\mathbb{Q}$, this inverse is again of the same form with rational entries.
    Thus we have shown that the subring is also closed under multiplicative inverses of nonzero elements.
    This is a field.
\end{proposition}

\begin{lemma}{
    For $n$ composite, $\zn$ is not a field if it has zero-divisors.
}
    If $\zn$ has a zero-divisor $[a]\neq [0]$, then there exists $[b]\neq [0]$ such that
    \[
        [a][b]=[0].
    \]
    If $\zn$ were a field, then $[a]$ would be a unit.
    Multiplying by $[a]^{-1}$ would give
    \[
        [b]=[0],
    \]
    a contradiction.
    Thus $\zn$ is not a field.
\end{lemma}

    {\Huge{From now on $\ri,\si$ is a ring and $\f$ is a field.}}\\

\begin{definition}[Integral Domain]
    Suppose $\ri$ is a commutative ring with $1\neq 0$.
    We say $\ri$ is an integral domain if $a\neq 0$ and $a\in\ri$ imply that $a$ is not a zero-divisor.
\end{definition}

We can think of these integral domain rings as being almost a field, but the only thing discerning them from being a field is the fact that not every nonzero element must have an inverse.
The only zero-divisor is $0_\ri$.
Remember that if there is a nonzero zero-divisor in $\ri$, then there is no way it can be a field, since all nonzero elements of a field must have an inverse, a.k.a.\ be a unit.

\begin{corollary}{
    $\f$ is an integral domain.
}
    Suppose $a,b\in\f$ with $ab=0$.
    Suppose $a\neq 0$.
    Then $a$ is a unit with inverse $a^{-1}$.
    Then
    \begin{align*}
        a^{-1}(ab) &= a^{-1}\cdot 0\\
        (a^{-1}a)b &= 0\\
        1b &= 0\\
        b &= 0.
    \end{align*}
    Thus if $ab=0$, then $a=0$ or $b=0$.
    So $\f$ is an integral domain.
\end{corollary}

Let us look into something called extensions.

\begin{definition}[Field Adjoins]
    We call something an adjoin given that suppose we have $\f=\mathbb{Q}$.
    Note this field is a subfield of $\mathbb{R}$.
    Then an extension of $\mathbb{Q}$ is taking an element of $\mathbb{R}\setminus\mathbb{Q}$ and adding it to $\mathbb{Q}$.
    An example of this is
    \[
        \mathbb{Q}\left[\sqrt{7}\right] \coloneqq \{a+b\sqrt{7}:a,b\in\mathbb{Q}\}.
    \]
\end{definition}

In fact, an exercise to do is to show that $\mathbb{Q}\left[\sqrt{7}\right]$ is a subfield.
Based on everything we have observed, we can say that $\zp$ is a field and $\zn$ is not even an integral domain when $n$ is composite.

\begin{axiom}[Pigeonhole Principle]
    If you have $n+1$ objects in $n$ slots, one slot will have more than $1$ element.
\end{axiom}

\begin{theorem}{
    A finite integral domain is a field.
}
    Let $F$ be a finite integral domain.
    We need to show that if $0\neq u\in F$, then $u$ has a multiplicative inverse.
    Consider the set
    \[
        \{u,u^2,u^3,\hdots\}.
    \]
    Suppose $F$ has $n$ elements.
    Then there must be repetition, so $u^k=u^m$ for some $m>k$.
    Thus
    \begin{align*}
        u^m-u^k &= 0\\
        u^k(u^{m-k}-1) &= 0.
    \end{align*}
    Since $F$ is an integral domain, then $u^k=0$ or $u^{m-k}-1=0$.
    Since $u\neq 0$, we must have $u^k\neq 0$.
    Then
    \begin{align*}
        u^{m-k}-1 &= 0\\
        u^{m-k} &= 1\\
        u(u^{m-k-1}) &= 1.
    \end{align*}
    Thus $u^{-1}=u^{m-k-1}$.
    So every nonzero element of $F$ has a multiplicative inverse.
    Therefore $F$ is a field.
\end{theorem}